\section{Introduction}
\label{sec:introduction}

A companion dissertation, \textit{Positive Governance}
\parencite{yuenyong2025positive}, uses structural equation modeling (SEM)
on World Values Survey data to identify a two-dimensional latent governance
space --- a State Capacity factor ($S$) and an Institutional Constraint
factor ($Z$, combining rule of law and democratic accountability) --- and
argues that this space, not a single one-dimensional governance ranking
such as the World Bank's Worldwide Governance Indicators, is the right unit
for comparative analysis. Within that space, the dissertation identifies
three qualitatively \textit{theorized} regions of empirically observed
governance configurations, which it labels Authoritarian Efficiency,
Constrained Liberalism, and Integrated Equilibrium. It is important, and a
central methodological commitment of this paper, to state precisely what
the dissertation itself claims for these three regions and what it does
not. The dissertation's own text is explicit that SEM ``does not classify
countries or predict trajectories; it defines the coordinate space,'' and
its own concluding chapter calls explicitly for \textit{future} work using
``unsupervised learning methods (e.g., cluster analysis, t-SNE) to
empirically identify governance archetypes.'' No clustering algorithm was
run on the SEM output to derive three trajectories; three is a qualitative,
interpretive typology laid over a statistically identified two-dimensional
space, not a statistically discovered partition count. We take this
distinction seriously throughout, because it is exactly the empirical
question a generative, agent-based search can speak to in a way a
cross-sectional SEM cannot: whether heterogeneous governance configurations,
and structured movement between them, can emerge \textit{endogenously} from
elementary micro-level mechanisms, with no typology, count, or macro-target
ever built in.

This paper reports a four-stage program of outcome-blind agent-based
search, structured around the dissertation's own governing equation,
$G = f(S, Z)$, extended by an unoperationalized external-influence term
$\beta$ the dissertation itself leaves open. We ask a bottom-up question the
dissertation's top-down SEM analysis cannot: if a population of
heterogeneous agents is given only elementary relational, ecological, and
organizational mechanisms --- cooperation, punishment, imitation, exit, a
resource ecology with real material accounting, proto-institutional rule
formation, delegation, monitoring, accountability, mobility, and
interpatch coordination --- and no encoding whatsoever of $S$, $Z$, or any
trajectory label, what persistent macroscopic governance patterns, if any,
emerge from the resulting dynamics, and does adding historically structured
pressure, sequencing, and entrenchment change the answer?

\subsection{From a synchronic null to a historically conditioned question}
\label{sec:intro-synchronic-to-historical}

The paper's first two milestones, Level 1.2 and Level 1.2b (in this
project's own internal numbering, continuing from an earlier development
sequence summarized in \cref{sec:model-genealogy}), tested a
\textit{synchronic} substrate: every simulated jurisdiction faces
materially the same menu of possible institutional configurations at
materially the same starting conditions, differing only in randomly drawn
traits and parameter draws, with no historically asymmetric pressure,
sequencing, or path-dependent lock-in represented anywhere. That search,
conducted twice --- once at a modest scale legitimately confounded by
near-universal ecological collapse, and once at roughly 24 times the scale
with that confound explicitly removed through exact ecological accounting
and ex-ante feasibility stratification --- did not independently recover a
small number of discrete, well-separated institutional forms. Both
milestones returned a \texttt{continuum} verdict, independently replicated
on fresh, disjoint data. We reported this as a positive, informative
finding, not a shortfall: the tested atomic governance ecology is
\textit{generatively sufficient} to produce institutions and heterogeneous
institutional outcomes from purely local rules, while showing, with real
statistical force, that generative sufficiency does not by itself imply
convergence onto a small number of recurring forms. We closed that
milestone by proposing, as the natural next hypothesis rather than a
finding already in hand, a formal decomposition

\begin{equation}
M_{\text{historical}} = M_{\text{general}} + H
\label{eq:historical-decomposition-intro}
\end{equation}

where $M_{\text{general}}$ is the synchronic substrate Level 1.2/1.2b
tested and $H$ denotes historically structured pressures, sequencing, and
entrenchment --- explicitly not implemented or tested at that time.

Level 1.3 and Level 1.3b, reported here for the first time, implement and
test $H$. They wire three pressure-transmission channels (security
pressure, bargaining pressure, and coordination-complexity pressure, each
via two of six candidate transmission mechanisms), institutional
entrenchment (age-dependent resistance to change), and dual
performance/procedural legitimacy ledgers into the same generative
substrate Level 1.2/1.2b already validated, and ask, generatively and then
causally, whether historically structured pressure and sequencing move the
institutional outcome space in ways a synchronic substrate cannot. Because
the observer layer's own $S$ and $Z$ composites ($S_{\text{obs}}$,
$Z_{\text{obs}}$; \cref{sec:observer-metrics-l13}) are constructed
independently of, and are not the same object as, the dissertation's SEM
latent factors, this program provides three complementary, non-duplicative
kinds of evidence about the dissertation's own $S$/$Z$ framework: whether a
generative mechanism set can produce heterogeneous configurations
distributed along axes that plausibly correspond to state capacity and
institutional constraint; whether a real, causal, historically structured
process shapes movement through that space; and whether the resulting
structure is discrete, in the way a three-trajectory typology would
require, or continuous.

\subsection{What this paper establishes, in sequence}
\label{sec:intro-sequence}

Six things are established across the paper's four linked milestones:

\begin{enumerate}
\item A general microscopic governance-generative agent-based model (ABM)
  was constructed, built strictly from agent-level mechanisms addressing
  cooperation, ecological viability, information, enforcement, delegation,
  accountability, mobility, and interpatch coordination
  (\cref{sec:model-methods}), with a structural firewall mechanically
  preventing any dissertation vocabulary, symbol, or trajectory label from
  ever entering the generative or observer layers.
\item An outcome-blind statistical search at Level 1.2/1.2b did not
  independently recover three discrete institutional types, or any small,
  stable number of them, in the synchronic substrate
  (\cref{sec:level12,sec:level12b}); both the initial and the ecologically
  stratified, 24-times-larger rediscovery returned an independently
  replicated \texttt{continuum} verdict.
\item Level 1.3 extended the substrate with historically structured
  pressure, sequencing, and entrenchment mechanisms and found, on pooled
  observational data, real ordering-dependent variation, a real
  (if construction-audited) $S_{\text{obs}}$--$Z_{\text{obs}}$ coupling, and
  a headline discovery result that itself failed independent fresh-data
  reproduction --- evidence leaning \texttt{continuum} again, but on a
  pooled, not purpose-designed, evidentiary base (\cref{sec:level13}).
\item Level 1.3b validated Level 1.3's findings with a purpose-designed,
  preregistered, eight-hypothesis program including, for the first time in
  this project, a sham-controlled causal intervention stage: erasing,
  reversing, or exposure-equalizing a run's realized pressure history
  produces real, persistent, non-decaying divergence from that run's own
  intact history, validated against a byte-identical sham control that
  itself introduces zero spurious divergence (\cref{sec:level13b}). History,
  in this precise causal sense, matters.
\item That same purpose-designed program, applying Level 1.2b's own frozen
  discreteness classifier unmodified, found the historically extended
  substrate's outcome space to be, if anything, \textit{more} decisively
  continuous than the synchronic substrate was: no tested $k$ from 1 to 8
  clears even the classifier's entry-level silhouette floor, cross-algorithm
  agreement is weak, and the leading eight-component partition does not
  reproduce across ecological strata (\cref{sec:level13b,sec:results}).
\item Taken together, the four milestones support a specific, bounded
  claim, stated precisely and repeatedly to avoid the very overclaiming this
  paper's own terminological discipline (\cref{sec:terminology-discipline}) is
  built to prevent: the governance space this model generates is
  historically structured, heterogeneous, and continuous --- not
  historically inert, and not discretely speciated into three, eight, or
  any other fixed count of institutional forms, using the precise
  vocabulary fixed in \cref{sec:terminology-discipline}
  (\cref{sec:discussion,sec:conclusion}).
\end{enumerate}

\subsection{Methodological triangulation, not duplication}
\label{sec:intro-triangulation}

We are explicit throughout that this paper's ABM evidence and the
dissertation's SEM evidence answer different questions at different
analytical levels, and that neither substitutes for the other.
\Cref{sec:disc-triangulation} develops this in full; in outline, the SEM
analysis \textit{identifies} a two-dimensional coordinate space from
cross-sectional survey data on real, already-formed states at one
historical moment; the ABM's generative layer establishes whether a
microscopic mechanism set is \textit{sufficient} to produce heterogeneous
configurations occupying an analogous space from nothing; and the ABM's
historical layer, uniquely, supports controlled causal experimentation ---
erasing or reversing a specific run's own history and comparing it against
itself, sham-validated --- that cross-sectional survey data structurally
cannot support at all. These three lines of evidence triangulate toward a
single account without duplicating each other, and no single one of them,
including this paper's own simulation results, is sufficient on its own to
settle whether real governance systems are, in fact, discretely speciated;
that stronger claim would require historically calibrated data this paper
does not use (\cref{sec:limitations}).

\subsection{Roadmap}
\label{sec:intro-roadmap}

Section~\ref{sec:theoretical-motivation} situates the project relative to
agent-based computational social science, common-pool-resource governance,
historical institutionalism and path dependence, and the methodological
literature on null results and discreteness testing in clustering.
Section~\ref{sec:model-genealogy} summarizes the level-by-level development
sequence, including Level 1.3/1.3b's own place in it.
Section~\ref{sec:model-methods} documents the full model and discovery
methodology, including the historical pressure/entrenchment/legitimacy
machinery and the extended eight-way decision-gate vocabulary.
Sections~\ref{sec:level12} and~\ref{sec:level12b} report the synchronic
Level 1.2/1.2b program; Sections~\ref{sec:level13}
and~\ref{sec:level13b} report the historical Level 1.3/1.3b program.
Section~\ref{sec:results} consolidates every headline quantitative result
across all four milestones. Section~\ref{sec:discussion} develops the
paper's theoretical synthesis: the richer generative state
$\mathcal{G}_t = [S_t, Z_t, H_t, E_t, L_t, P_t, R_t, C_t, \ldots]$, the
triangulation logic linking SEM identification, ABM generative sufficiency,
and historical causal experimentation, and a precise, bounded reading of
the relationship to the dissertation's own three trajectories.
Sections~\ref{sec:limitations}--\ref{sec:reproducibility} close with
limitations, conclusions, and reproducibility information. Appendices
provide a compact ODD-style model summary extended for Level 1.3/1.3b, full
parameter and outcome-count tables, the recurrence and decision-gate
procedures, the ledger accounting identities, and a condensed evidence
ledger mapping every numerical claim in this paper to its source file and
field; the complete claim-by-claim manifest is maintained separately in
\texttt{claim\_evidence\_manifest.csv}.
